% Options for packages loaded elsewhere
\PassOptionsToPackage{unicode}{hyperref}
\PassOptionsToPackage{hyphens}{url}
\PassOptionsToPackage{dvipsnames,svgnames*,x11names*}{xcolor}
%
\documentclass[
  11pt,
]{article}
\usepackage{lmodern}
\usepackage{amssymb,amsmath}
\usepackage{ifxetex,ifluatex}
\ifnum 0\ifxetex 1\fi\ifluatex 1\fi=0 % if pdftex
  \usepackage[T1]{fontenc}
  \usepackage[utf8]{inputenc}
  \usepackage{textcomp} % provide euro and other symbols
\else % if luatex or xetex
  \usepackage{unicode-math}
  \defaultfontfeatures{Scale=MatchLowercase}
  \defaultfontfeatures[\rmfamily]{Ligatures=TeX,Scale=1}
\fi
% Use upquote if available, for straight quotes in verbatim environments
\IfFileExists{upquote.sty}{\usepackage{upquote}}{}
\IfFileExists{microtype.sty}{% use microtype if available
  \usepackage[]{microtype}
  \UseMicrotypeSet[protrusion]{basicmath} % disable protrusion for tt fonts
}{}
\makeatletter
\@ifundefined{KOMAClassName}{% if non-KOMA class
  \IfFileExists{parskip.sty}{%
    \usepackage{parskip}
  }{% else
    \setlength{\parindent}{0pt}
    \setlength{\parskip}{6pt plus 2pt minus 1pt}}
}{% if KOMA class
  \KOMAoptions{parskip=half}}
\makeatother
\usepackage{xcolor}
\IfFileExists{xurl.sty}{\usepackage{xurl}}{} % add URL line breaks if available
\IfFileExists{bookmark.sty}{\usepackage{bookmark}}{
\usepackage{hyperref}
}
\hypersetup{
  pdftitle={TP table de données, requêtes et tri},
  pdfauthor={Première NSI Lycée du Parc},
  colorlinks=true,
  linkcolor=Maroon,
  filecolor=Maroon,
  citecolor=Blue,
  urlcolor=Blue,
  pdfcreator={LaTeX via pandoc}}
\urlstyle{same} % disable monospaced font for URLs
\usepackage[top=20mm,left=20mm,right=20mm,heightrounded]{geometry}
\usepackage{listings}
\newcommand{\passthrough}[1]{#1}
\lstset{defaultdialect=[5.3]Lua}
\lstset{defaultdialect=[x86masm]Assembler}
\setlength{\emergencystretch}{3em} % prevent overfull lines
\providecommand{\tightlist}{%
  \setlength{\itemsep}{0pt}\setlength{\parskip}{0pt}}
\setcounter{secnumdepth}{5}

\title{TP table de données, requêtes et tri}
\usepackage{etoolbox}
\makeatletter
\providecommand{\subtitle}[1]{% add subtitle to \maketitle
  \apptocmd{\@title}{\par {\large #1 \par}}{}{}
}
\makeatother
\subtitle{Traitement de données en tables}
\author{Première NSI Lycée du Parc}
\date{}

%%%jolis boites

\usepackage{fancybox, graphicx}



%%%%%%%%%%%%%%%%Packages et Macros Frederic%%%%%%%%%%%%%%%%%%%%%%%%%%%%%


%%%%Insertion de liens hypertextes %%%%

            
%%%%%%%%%%PSTricks%%%%%%%%%%%%

\usepackage{pstricks,pst-plot,pst-text,pst-tree,pst-eps,pst-fill,pst-node,pst-math,pstricks-add,pst-xkey,pst-eucl}


%%%%%%%Tikz%%%%%%%%%%%%%%%
\usepackage{pgf,tikz,tkz-tab}
% Pour les tableaux de signes ou de variations avec tkz-tab voir https://zestedesavoir.com/tutoriels/439/des-tableaux-de-variations-et-de-signes-avec-latex/#1-13389_tikz-un-package-qui-en-a-dans-le-ventre
\usetikzlibrary{arrows}
\usetikzlibrary{shapes.geometric}
\usetikzlibrary{shapes.geometric}
\usetikzlibrary{petri}
\usetikzlibrary{decorations}
\usetikzlibrary{arrows}
\usetikzlibrary{math}
 %Variables must be declared in a tikzmath environment but
       % can be used outside
%       \tikzmath{int \n; \n = 508; \x1 = 1; \y1 =1; 
%                   %computations are also possible
%                    \x2 = \x1 + 1; \y2 =\y1 +3; } 


%%%%%%%%%%%%%%%%%%%%%%%%%%%%%%%%%%%%%%%%
%%%%%%%%%%%Commandes Tikz Perso%%%%%%%%%%%%%%%

% Définition des nouvelles options xmin, xmax, ymin, ymax
% Valeurs par défaut : -3, 3, -3, 3
\tikzset{
xmin/.store in=\xmin, xmin/.default=-3, xmin=-3,
xmax/.store in=\xmax, xmax/.default=3, xmax=3,
ymin/.store in=\ymin, ymin/.default=-3, ymin=-3,
ymax/.store in=\ymax, ymax/.default=3, ymax=3,
}
% Commande qui trace la grille entre (xmin,ymin) et (xmax,ymax)
\newcommand {\grille}[2]
{\draw[help lines,black, thick] (\xmin,\ymin) grid[xstep=#1, ystep=#2] (\xmax,\ymax);}
% Commande \axes
\newcommand {\axes} {
\draw[->,very thick] (\xmin,0) -- (\xmax,0);
\draw[->,very thick] (0,\ymin) -- (0,\ymax);
\draw (0.95*\xmax, 0) node[above] {};
\draw (0, 0.95*\ymax) node[left] {};
}
% Commande qui limite l?affichage à (xmin,ymin) et (xmax,ymax)
\newcommand {\fenetre}
{\clip (\xmin,\ymin) rectangle (\xmax,\ymax);}

%Exemple d'utilisation

%\begin{center}
%\begin{tikzpicture} [xmin=-2,xmax=2,ymin=0,ymax=5]
%\grille{1} \axes \fenetre
%\draw plot[smooth] (\x,\x^2);
%\end{tikzpicture}
%\end{center}

%style pour la perspective cavalière française
%voir Tikz pour l'impatient page 68
\tikzset{math3d/.style=
{x= {(-0.353cm,-0.353cm)}, z={(0cm,1cm)},y={(1cm,0cm)}}}

%%%%%%%Symbole pour code calculatrice%%%%%%

%Flèche remplie pour défilement de menu

\newcommand{\flechefillright}{
\begin{tikzpicture}[scale=0.15] \fill (0,0)--(2,1)--(0,2)--cycle;
\end{tikzpicture}}

%%%%%%%%%%%%%Symboles pour calculatrice Casio%%%%
\newcommand{\execasio}{\Pisymbol{psy}{191}} %Retour chariot
\newcommand{\dispcasio}{\begin{pspicture}(.1,.1)\pspolygon*(.1,0)(.1,.1)\end{pspicture}} %Triangle « Disp »
\newcommand{\dispcasiotikz}{\begin{tikzpicture}[scale=0.2]
\fill (0,0) -- (1,0) -- (1,1) -- cycle;
\end{tikzpicture}} %Triangle « Disp »
%

%Fleche entre deux lignes, d'apres 'un bon petit' : http://forum.mathematex.net/latex-f6/fleches-entre-deux-lignes-pour-resolution-d-equation-t10283.html#p99817
\newcommand\addnode[1]{\Rnode{#1}{}}
\newcommand\linknode[3]{\ncbar[angleA=0,angleB=0,nodesep=1ex,arm=10ex,offset=-2pt]{->}{#1}{#2}\Aput{\vphantom{x}#3}}


%%Commande pour touche de calculatrice

\newcommand\tc[1]{%
{
\begin{tikzpicture}
\node[draw,rectangle,rounded corners=3pt] (P) at (0,0){#1};
\end{tikzpicture}
}
}

%%%%%%%%%%%%%%%%%%%%%%%%%%%%%%%%%%%%%%%%
%%%%%%%%%%%Fin Commandes Tikz%%%%%%%%%%%%%%%


%%%%%%%%%%%%Specifiques%%%%%%%%%%%
\usepackage{wrapfig}
%pour insérer une figure à droite ou à gauche d'un texte
%\begin{wrapfigure}[nb lignes]{placement l,r,c,i(inside),o(outside)}[overhang]{width}
%ce package fonctionne mal à proximité des listes
%%%%%%%%%%%%%%%%%%%%%%%%%%%%%%%%%%%%%

%%%%%Environnements et symboles spéciaux pour faire joli%%%%%%

%%%Bclogo, pour des environnements + jolis avec insertion de logo%%%%
%Dépendances de  bclogo
\usepackage{xkeyval}  
\usepackage{etoolbox}
\usepackage{ifpdf}
\usepackage[framemethod=tikz]{mdframed}
\usepackage[tikz]{bclogo}

%\newcommand\bcpython{\includegraphics[width=17pt]{/home/fjunier/Maths/python-logo.eps}}
\newcommand\bcpython{\includegraphics[width=17pt]{/home/fjunier/Maths/python-logo.png}}
%\newcommand\bcpython{\includegraphics[width=17pt]{/home/frederic/Maths/python-logo.png}}

%% Framed
\usepackage{framed}  %Le package « framed» Crée 3 nouveaux environnements, qui se comportent comme des minipage de largeur \linewidth, mais permettant en plus de se casser entre plusieurs pages.     * framed : avec un cadre autour;     * shaded : avec un fonc coloré (il faut définir la couleur shadecolor);     * leftbar : avec une barre le long du côté gauche.

%%%%%%%%%%%%%%%%%%%Présentation de codes sources%%%%%%%%%%%%%%%%%
\usepackage{listings}
%On utilise l?environnement lstlisting pour insérer
%un code source.
%En plus de l?environnement lstlisting, on peut également utiliser la
%commande \lstinline qui fonctionne comme la commande \verb, en ce
%sens qu?on peut utiliser n?importe quel caractère comme délimiteur. Enfin,
%la commande \lstinputlisting permet de charger un code source depuis
%un fichier externe.
%Il y a deux manières de préciser des options : soit via l?option de l?envi-
%ronnement ou de la commande, soit en utilisant la commande \lstset
%qui permet de définir des options de manière globale.

\lstset{ %
  language=Python,                % the language of the code
  basicstyle=\ttfamily,           % the size of the fonts that are used for the code
  %numbers=left,                   % where to put the line-numbers
  numberstyle=\tiny,  % the style that is used for the line-numbers
  %stepnumber=2,                   % the step between two line-numbers. If it's 1, each line 
                                  % will be numbered
  %numbersep=5pt,                  % how far the line-numbers are from the code
  backgroundcolor=\color{white},      % choose the background color. You must add \usepackage{color}
  showspaces=false,               % show spaces adding particular underscores
  showstringspaces=false,         % underline spaces within strings
  showtabs=false,                 % show tabs within strings adding particular underscores
  frame=single,                   % adds a frame around the code
  rulecolor=\color{black},        % if not set, the frame-color may be changed on line-breaks within not-black text (e.g. comments (green here))
  tabsize=4,                      % sets default tabsize to 2 spaces
  captionpos=b,                   % sets the caption-position to bottom
  breaklines=true,                % sets automatic line breaking
  breakatwhitespace=false,        % sets if automatic breaks should only happen at whitespace
  %title=\lstname,                   % show the filename of files included with \lstinputlisting;
                                  % also try caption instead of title
  breakindent=1cm,
  keywordstyle=\color{blue},          % keyword style
  commentstyle=\color{red},       % comment style
  %stringstyle=\ttfamily\color{green},         % string literal style
  escapeinside={\%*}{*)},            % if you want to add LaTeX within your code
  morekeywords={*,...},              % if you want to add more keywords to the set
  deletekeywords={...}              % if you want to delete keywords from the given language
  upquote=true,columns=flexible,
xleftmargin=1cm,xrightmargin=1cm,
 inputencoding=utf8,			%Les lignes qui suivent sont pour le codage utf8
  extendedchars=true,
  literate=%
            {é}{{\'{e}}}1
            {è}{{\`{e}}}1
            {ê}{{\^{e}}}1
            {ë}{{\¨{e}}}1
            {û}{{\^{u}}}1
            {ù}{{\`{u}}}1
            {â}{{\^{a}}}1
            {à}{{\`{a} }}1
            {î}{{\^{i}}}1
            {ô}{{\^{o}}}1
            {ç}{{\c{c}}}1
            {Ç}{{\c{C}}}1
            {É}{{\'{E}}}1
            {Ê}{{\^{E}}}1
            {À}{{\`{A}}}1
            {Â}{{\^{A}}}1
            {Î}{{\^{I}}}1
}

\lstdefinestyle{rond}{
  numbers=none,
  backgroundcolor=\color{gristclair},
  frameround =tttt
}

\lstdefinestyle{compil}{
  numbers=none,
  backgroundcolor=\color{gristclair}
}
%\lstset{language=Python,basicstyle=\small , frame=single,tabsize=4,showspaces=false,showtabs=false,showstringspaces=false,numbers=left,numberstyle=\tiny , extendedchars=true}



%%%%%%%%%%%%%%%%%%%%%%%%%%%%%%%%%%%%%%%%%%%%%%%%%%%%%%%%%%%%%%%%%%%%%%%%
%%%%%%%%%%%%%%%%%%%%Environnements persos%%%%%%%%%%%%%%%%%%%%%%%%%%%%%%%%
%Syntaxe :
%\newenvironment{nom}[nombre d'args][defaut]{definitions initiales}{definitions finales}
%definitions intiales sont les commandes appelées par \begin{nom}
%Definitions finales sont les commandes appelées par \end{nom}

%%%%%%%%%%%%%%%%Définitions des environnemts persostheoreme, exemple ..%%%%
%%%% Exercice avec encadré %%%%
\newcounter{exo}
\newenvironment{exercice}[1]
{\par \medskip   \addtocounter{exo}{1} \noindent  
\begin{bclogo}[arrondi =0.1,   noborder = true, logo=\bccrayon, marge=4]{~\textbf{Exercice} \textbf{\theexo} {\itshape #1} }  \par}
{
\end{bclogo}
 \par \bigskip }

%%Axiomes, Theoremes, Propriété, Définition, Methode, Preuve


\newenvironment{axiome}[1]
{\par \medskip   \begin{leftbar} \noindent \underline{\textbf{Axiome}}\hspace{0.5cm}{\itshape #1}   \vspace*{10pt} \par }
{\end{leftbar}  \par \medskip }


\newcounter{thme}
\newenvironment{theoreme}[1]
{\par \medskip  \addtocounter{thme}{1} \noindent  
\begin{bclogo}[arrondi =0.1,  ombre = true, barre=none, logo=\bcbook, marge=4]{~\textbf{Théorème} \textbf{\thethme} {\itshape #1} }   \par}
{
\end{bclogo}
 \par \bigskip}

 \newenvironment{theoremedef}[1]
{\par \medskip   \addtocounter{thme}{1} \noindent  
\begin{bclogo}[arrondi =0.1,  ombre = true, barre=none, logo=\bcbook, marge=4]{~\textbf{Théorème-Définition} \textbf{\thethme} {\itshape #1} }   \par}
{
\end{bclogo}
 \par \bigskip }
 
\newcounter{prop}
\newenvironment{propriete}[1]
{\par \medskip   \addtocounter{prop}{1} \noindent  
\begin{bclogo}[arrondi =0.1,  ombre = true, barre=none, logo=\bcbook, marge=4]{~\textbf{Propriété} \textbf{\theprop} {\itshape #1} }   \par}
{
\end{bclogo}
 \par \bigskip }


\newenvironment{corollaire}[1]
{\par \medskip   \noindent  
\begin{bclogo}[arrondi =0.1,  ombre = true, barre=none, logo=\bcbook, marge=4]{~\textbf{Corollaire} {\itshape #1} } \par }
{
\end{bclogo}
 \par \bigskip }

\newenvironment{demo}[1]
{\par \medskip   \noindent  
\begin{bclogo}[arrondi =0.1,  ombre = true, barre=zigzag, noborder = true, logo=\bcloupe, marge=0]{~\textbf{Démonstration} {\itshape #1} } \par \vspace{10pt}}
{
\end{bclogo}
 \par \bigskip }

\newcounter{activite}
\newenvironment{activite}[1]
{\par \medskip   \noindent   \addtocounter{activite}{1}
\begin{bclogo}[arrondi =0.1,   noborder = true, logo=\bcvelo, marge=4]{~\textbf{Activité} \textbf{\theactivite} {\itshape #1} }  \par}
{
\end{bclogo}
 \par \bigskip }


\newenvironment{synthese}
{\par \medskip   \noindent   
\begin{bclogo}[arrondi =0.1,   noborder = true, logo=\bccle, marge=4]{~\textbf{Synthèse}   }  \par}
{
\end{bclogo}
 \par \bigskip }
 
 
\newcounter{rque}
\newenvironment{remarque}
{\par \medskip    \addtocounter{rque}{1} \noindent  
\begin{bclogo}[arrondi =0.1,  ombre = true, barre=snake, noborder = true, logo=\bcinfo, marge=0]{~\textbf{Remarque} \textbf{\therque}}  \par }
{
\end{bclogo}
 \par \bigskip }

\newcounter{def}
\newenvironment{definition}[1]
{\par \medskip   \addtocounter{def}{1} \noindent  
\begin{bclogo}[arrondi =0.1,  ombre = true, barre=none, logo=\bcbook, marge=4]{~\textbf{Définition} \textbf{\thedef} {\itshape #1} }  \par}
{
\end{bclogo}
 \par \bigskip }
 
 
 \newcounter{cours}
\newenvironment{cours}[1]
{\par \medskip   \addtocounter{cours}{1} \noindent  
\begin{bclogo}[arrondi =0.1,  ombre = true, barre=none, logo=\bcbook, marge=4]{~\textbf{Point de cours} \textbf{\thecours} {\itshape #1} }  \par}
{
\end{bclogo}
 \par \bigskip }
 
 

\newenvironment{introduction}
{\par \medskip    \noindent  
 \begin {bclogo}[couleur = blue!5 , arrondi =0.1,logo=\bcrosevents, marge=4] {~\textbf{Introduction}    }
 \par }
{
\end{bclogo}
 \par \bigskip }
 
\newenvironment{memo}[1]
{\par \medskip    \noindent  
\begin{bclogo}[arrondi =0.1,  ombre = true, barre=none, logo=\bccle, marge=4]{~\textbf{À retenir}  {\itshape #1} }  \par}
{
\end{bclogo}
 \par \bigskip }
 
\newcounter{exple}
\newenvironment{exemple}[1]
{\par \medskip   \addtocounter{exple}{1} \noindent  
\begin{bclogo}[arrondi =0.1,   noborder = true, logo=\bclampe, marge=4]{~\textbf{Exemple} \textbf{\theexple} {\itshape #1} }  \par}
{
\end{bclogo}
 \par \bigskip }




\newcounter{alg}
\newenvironment{algorithme}[1]
{\par \medskip   \addtocounter{alg}{1} \noindent  
 \begin {bclogo}[noborder = true, barre=zigzag,logo=\bcpython, marge=4] {~\textbf{Algorithmique} \textbf{\thealg} {\itshape #1} }  \par}
{
\end{bclogo}
 \par \bigskip }

\newcounter{prog}
\newenvironment{programme}[1]
{\par \medskip   \addtocounter{prog}{1} \noindent  
 \begin {bclogo}[noborder = true, barre=zigzag,logo=\bcpython, marge=4] {~\textbf{Programme} \textbf{\theprog} {\itshape #1} }  \par  \bigskip}
{
\end{bclogo}
 \par \bigskip }
 
\newcounter{logi}
\newenvironment{logique}[1]
{\par \medskip   \addtocounter{logi}{1} \noindent  
 \begin {bclogo}[noborder = true, barre=zigzag,logo=\bclampe, marge=4] {~\textbf{Logique} \textbf{\thelogi} {\itshape #1} }  \par}
{
\end{bclogo}
 \par \bigskip }


\newenvironment{methode}[1]
{\par \medskip    \noindent  
 \begin {bclogo}[arrondi =0.1,logo=\bcoutil, marge=4,noborder = true] {~\textbf{Méthode}   {\itshape #1} }  \par}
{
\end{bclogo}
 \par \bigskip }


\newcounter{histo}
\newenvironment{histoire}[1]
{\par \medskip   \addtocounter{histo}{1} \noindent  
 \begin {bclogo}[couleur = blue!10 , arrondi =0.1,logo=\bchorloge, marge=4] {~\textbf{Histoire} \textbf{\thehisto} {\itshape #1} }  \par}
{
\end{bclogo}
 \par \bigskip }




%Environnement contenu pour un document présentant une progression annuelle
\newenvironment{contenu}
{\par \medskip   \begin {bclogo}[ noborder = true,logo=\bccrayon] \noindent {\large \textbf{Contenu de la séance}} \vspace*{10pt} \par  }
{\end{bclogo}  \par \medskip }

%Environnement programme pour un document présentant une progression annuelle
%\newenvironment{programme}
%{\par \medskip   \begin {bclogo}[ noborder = true, barre=zigzag,logo=\bcinfo] \noindent {\large \textbf{Programme officiel}} \vspace*{10pt} \par  }
%{\end{bclogo}  \par \medskip }

%Environnement programme pour un document présentant une progression annuelle
\newenvironment{ressource}
{\par \medskip   \begin {bclogo}[ noborder = true,logo=\bcbook] \noindent {\large \textbf{Ressources}}\\vspace*{10pt} \par }
{\end{bclogo}  \par \medskip }




%%%%%%%%%%%%%%%%%%Maths divers%%%%%%%%%%%%%%%%%%%%%%%%%
%%%%%%%%%%%%%Nombres%%%%%%%%%%%%%%%%

%Ensemble prive de...
%\newcommand{\prive}{\boi}%{\backslash}

%Ensembles de nombres%%%%%%%%%%%%%%%%%
\newcommand{\R}{\mathbb{R}}
\newcommand{\N}{\mathbb{N}}
\newcommand{\D}{\mathbb{D}}
\newcommand{\Z}{\mathbb{Z}}
\newcommand{\Q}{\mathbb{Q}}
%\newcommand{\C}{\mathbb{C}}
\newcommand{\df}{~\ensuremath{]0;+\infty[}~}
\newcommand{\K}{\mathbb{K}}

%%%%%%%%Arithmetique%%%%%%%%%%
%PGCD, PPCM
\newcommand{\PGCD}{\mathop{\rm PGCD}\nolimits}
\newcommand{\PPCM}{\mathop{\rm PPCM}\nolimits}

%Intervalles
\newcommand{\interoo}[2]{]#1\, ;\, #2[}
\newcommand{\Interoo}[2]{\left]#1\, ;\, #2\right[}
\newcommand{\interof}[2]{]#1\, ;\, #2]}
\newcommand{\Interof}[2]{\left]#1\, ;\, #2\right]}
\newcommand{\interfo}[2]{[#1\, ;\, #2[}
\newcommand{\Interfo}[2]{\left[#1\, ;\, #2\right[}
\newcommand{\interff}[2]{[#1\, ;\, #2]}
\newcommand{\Interff}[2]{\left[#1\, ;\, #2\right]}
%\newcommand\interentiers #1#2{[\! [#1\, ;\, #2]\! ]}
\newcommand{\interentiers}[2]{\llbracket #1\, ;\, #2\rrbracket}
%


%%%%%%%%%%%%%%Nombres complexes%%%%%

\newcommand{\ic}{\text{i}}
%\newcommand{\I}{\text{i}}
\newcommand{\im}[1]{\text{Im}\left(#1\right)}
\newcommand{\re}[1]{\text{Re}\left(#1\right)}
\newcommand{\Arg}[1]{\text{arg}\left(#1\right)}
\newcommand{\Mod}[1]{\left[#1\right]}
%Parties entière, réelle, imaginaire, nombre i
\newcommand{\ent}[1]{\text{E}\left(#1\right)}
\renewcommand{\Re}{\mathop{\rm Re}\nolimits}
\renewcommand{\Im}{\mathop{\rm Im}\nolimits}
\renewcommand{\i}{\textrm{i}}

%%%%%%%%%%%Probabilites et statistiques%%%%%
\newcommand{\loibinom}[2]{\mathcal{B}\left(#1\ ; \ #2 \right)}
\newcommand{\loinorm}[2]{\mathcal{N}\left(#1\ ; \ #2 \right)}
\newcommand{\loiexp}[1]{\mathcal{E}\left(#1\right)}
\newcommand{\proba}[1]{\mathbb{P}\big(#1\big)}
\newcommand{\probacond}[2]{\mathbb{P}_{#2}\big(#1\big)}
\newcommand{\esperance}[1]{\mathbb{E}\left(#1\right)}
\newcommand{\variance}[1]{\mathbb{V}\left(#1\right)}
\newcommand{\ecart}[1]{\sigma\left(#1\right)}
\newcommand{\dnormx}{\frac{1}{\sqrt{2\pi}} \text{e}^{-\frac{x^2}{2}}}
\newcommand{\dnormt}{\frac{1}{\sqrt{2\pi}} \text{e}^{-\frac{t^2}{2}}}

%Covariance
\newcommand{\cov}{\mathop{\rm cov}\nolimits}
%


%%%%%%%%%%Analyse%%%%%%%%%%%

%%%%%%%%%%%Courbe%%%%%%%%%%%%
\newcommand{\courbe}[1]{\ensuremath{\mathcal{C}_{#1}}}

%%%%%%%Fonction exponentielle%%%%%
\newcommand{\fe}{~fonction exponentielle~}
\newcommand{\e}{\text{e}}

%Fonction cotangente
\newcommand{\cotan}{\mathop{\rm cotan}\nolimits}
%%%%%%%%%%%%%%%%%%%%%%%%%%%%%%%%%%%%%%%%%
%
%Fonctions hyperboliques
\newcommand{\ch}{\mathop{\rm ch}\nolimits}
\newcommand{\sh}{\mathop{\rm sh}\nolimits}


%%%%%%%%%%%%%%Limites%%%%%%
\newcommand{\limite}[2]{\lim\limits
_{x \to #1} #2}
\newcommand{\limitesuite}[1]{\lim\limits
_{n \to +\infty} #1}
\newcommand{\limiteg}[2]{\lim\limits
_{\substack{x \to #1 \\ x < #1 }} #2}
\newcommand{\limited}[2]{\lim\limits
_{\substack{x \to #1 \\ x > #1 }} #2}

%%%%%%%%%%Continuité%%%%%%%%%%%
\newcommand{\TVI}{théorème des valeurs intermédiaires}

%%%%%%%%%%%Suites%%%%%%%%%%%%
\newcommand{\suite}[1]{\ensuremath{\left(#1_{n}\right)}}
\newcommand{\Suite}[2]{\ensuremath{\left(#1\right)_{#2}}}
%

%%%%%%%%%%%%%%%Calcul intégral%%%%%%
\newcommand{\dx}{\ensuremath{\text{d}x}}		% dx
\newcommand{\dt}{\ensuremath{\text{d}t}}		% dt
\newcommand{\dtheta}{\ensuremath{\text{d}\theta}}		% dtheta
\newcommand{\dy}{\ensuremath{\text{d}y}}		% dy
\newcommand{\dq}{\ensuremath{\text{d}q}}		% dq

%%%Intégrale%%%
\newcommand{\integralex}[3]{\int_{#1}^{#2} #3 \ \dx}
\newcommand{\integralet}[3]{\int_{#1}^{#2} #3 \ \dt}
\newcommand{\integraletheta}[3]{\int_{#1}^{#2} #3 \ \dtheta}

%%%%%Equivalent%%
\newcommand{\equivalent}[1]{\build\sim_{#1}^{}}

%o et O%%%%
\renewcommand{\o}[2]{\build o_{#1\to #2}^{}}
\renewcommand{\O}[2]{\build O_{#1\to #2}^{}}



%%%%%%%%%%%%%%%Geometrie%%%%%%%%%%%%%%%%%%%%%%%

%%%%%%%%%%%%%%%Reperes%%%%%%%%%%%%%%
\def\Oij{\ensuremath{\left(\text{O},~\vect{\imath},~\vect{\jmath}\right)}}
\def\Oijk{\ensuremath{\left(\text{O},~\vect{\imath},~ \vect{\jmath},~ \vect{k}\right)}}
\def\Ouv{\ensuremath{\left(\text{O},~\vect{u},~\vect{v}\right)}}
\renewcommand{\ij}{(\vec\imath\, ;\vec\jmath\,)}
\newcommand{\ijk}{(\vec\imath\, ;\vec\jmath\, ;\vec k\,)}
\newcommand{\OIJ}{(O\,;\, I\,;\, J\,)}
\newcommand{\repere}[3]{\big(#1\, ;\,\vect{#2} ;\vect{#3}\big)}
\newcommand{\reperesp}[4]{\big(#1\, ;\,\vect{#2} ;\vect{#3} ;\vect{#4}\big)}

%%%%%%%%%Coordonnees%%%%%%%%%%%%%%
\newcommand{\coord}[2]{(#1\, ;\, #2)}
\newcommand{\bigcoord}[2]{\big(#1\, ;\, #2\big)}
\newcommand{\Coord}[2]{\left(#1\, ;\, #2\right)}
\newcommand{\coordesp}[3]{(#1\, ;\, #2\, ;\, #3)}
\newcommand{\bigcoordesp}[3]{\big(#1\, ;\, #2\, ;\, #3\big)}
\newcommand{\Coordesp}[3]{\left(#1\, ;\, #2\, ;\, #3\right)}
\newcommand{\Vcoord}[3]{\begin{pmatrix} #1 \\ #2 \\ #3 \end{pmatrix}}
%Symboles entre droites
%\newcommand{\paral}{\sslash}
\newcommand{\paral}{\mathop{/\!\! /}}
%

%%%%%%%%%Produit scalaire, Angles%%%%%%%%%%
\newcommand{\scal}[2]{\vect{#1} \, \cdot \, \vect{#2}}
\newcommand{\Angle}[2]{\left(\vect{#1} \, , \, \vect{#2}\right)}
\newcommand{\Anglegeo}[2]{\left(\widehat{\vect{#1} \, ; \, \vect{#2}}\right)}
\renewcommand{\angle}[1]{\widehat{#1}}
\newcommand{\anglevec}[2]{\left(\vect {#1}\, ,\,\vect {#2} \right)}
\newcommand{\anglevecteur}[2]{(#1\, , \, #2)}
\newcommand{\Anglevec}[2]{(\vecteur{#1}\, ,\,\vecteur{#2})}
\newcommand{\prodscal}[2]{#1 \, \cdot \, #2}
%


%Arc
%\newcommand{\arc}[1]{\wideparen{#1}}
\newcommand{\arcoriente}[1]{\overset{\curvearrowright}{#1}}
%
%


%%%%%%%%%%%%%%%Normes%%%%%%%%%%%%%%%%
\newcommand{\norme}[1]{\left\| #1\right\|}
\newcommand{\normebis}[1]{\delim{2pt}{\|}{9pt}\! #1\delim{2pt}{\|}{9pt}}
\newcommand{\normetriple}[1]{\left |\kern -.07em\left\| #1\right |\kern -.07em\right\|}
\newcommand{\valabs}[1]{\big| \, #1 \, \big|}
%

%%%%%%%%%%%%%%%%%%%%%%%%%%%Degré%%%%%%
%\newcommand{\Degre}{\ensuremath{^\circ}}
%La commande \degre est déjà définie dans le package babel

%%%%%%%%%%Vecteurs%%%%%%%%%%%
\newcommand{\vect}[1]{\mathchoice%
{\overrightarrow{\displaystyle\mathstrut#1\,\,}}%
{\overrightarrow{\textstyle\mathstrut#1\,\,}}%
{\overrightarrow{\scriptstyle\mathstrut#1\,\,}}%
{\overrightarrow{\scriptscriptstyle\mathstrut#1\,\,}}}



%%%%%%%%%%%%%Algebre%%%%%%%%%%%%%%%


%%%%%%%%%%Systemes%%%%%%%%%%%
%Systemes
\newcommand{\sys}[2]{
\left\lbrace
 \begin{array}{l}
  \negthickspace\negthickspace #1\\
  \negthickspace\negthickspace #2\\
 \end{array}
\right.\negthickspace\negthickspace}
\newcommand{\Sys}[3]{
\left\lbrace
 \begin{array}{l}
  #1\\
  #2\\
  #3\\
 \end{array}
\right.}
\newcommand{\Sysq}[4]{
\left\lbrace
 \begin{array}{l}
  #1\\
  #2\\
  #3\\
  #4\\
 \end{array}
\right.}
%
%

%%%%%%%%%%%%%%%%Matrices%%%%%%%%%%%%%%%%%%
%Comatrice
\newcommand{\com}{\mathop{\rm com}\nolimits}
%
%
%Trace
\newcommand{\tr}{\mathop{\rm tr}\nolimits}
%
%
%Transposee
\newcommand{\transposee}[1]{{\vphantom{#1}}^t\negmedspace #1}
%
%
%Noyau
\newcommand{\Ker}{\mathop{\rm Ker}\nolimits}
%
%

%
%Matrices
\newcommand{\Mn}{\mathcal M_n}
\newcommand{\matrice}[4]{
\left(
 \begin{array}{cc}
  #1 & #2 \\
  #3 & #4
 \end{array}
\right)}

\newcommand{\Matrice}[9]{
\left(
 \begin{array}{ccc}
  #1 & #2 & #3\\
  #4 & #5 & #6\\
  #7 & #8 & #9
 \end{array}
\right)}
\newcommand{\Vect}[3]{
\left(\negmedspace
 \begin{array}{c}
  #1\\
  #2\\
  #3
 \end{array}\negmedspace
\right)}
\newcommand{\Ideux}{\matrice{1}{0}{0}{1}}
\newcommand{\Itrois}{\Matrice{1}{0}{0}{0}{1}{0}{0}{0}{1}}
%
%
%Determinants
\newcommand{\determinant}[4]{
\left|
 \begin{array}{cc}
  #1 & #2 \\
  #3 & #4
 \end{array}
\right|}
\newcommand{\Determinant}[9]{
\left|
 \begin{array}{ccc}
  #1 & #2 & #3\\
  #4 & #5 & #6\\
  #7 & #8 & #9
 \end{array}
\right|}

\begin{document}
\maketitle

\renewcommand*\contentsname{Table des matières}
{
\hypersetup{linkcolor=}
\setcounter{tocdepth}{3}
\tableofcontents
}
\hypertarget{cruxe9dits}{%
\section*{Crédits}\label{cruxe9dits}}
\addcontentsline{toc}{section}{Crédits}

\emph{Ce TP est largement inspiré :}

\begin{itemize}
\tightlist
\item
  \emph{des chapitres 16 et 17 du manuel NSI de la collection Tortue
  chez Ellipse, auteurs : Ballabonski, Conchon, Filliatre, N'Guyen ;}
\item
  \emph{du document d'accompagnement des programmes de NSI disponible
  sur Eduscol}
\end{itemize}

\hypertarget{requuxeates-dans-une-table-de-donnuxe9es}{%
\section{Requêtes dans une table de
données}\label{requuxeates-dans-une-table-de-donnuxe9es}}

\begin{definition}{}

Dans le cours, on a vu comment charger une \textbf{table de données},
enregistrée dans un \textbf{fichier CSV}, dans une structure de données
permettant un accès efficace aux \textbf{enregistrements} et aux valeurs
de leurs \textbf{attributs}. En Python, une structure de données bien
adaptée est un \textbf{tableau de dictionnaires}. Les opérations sur les
tableaux et les dictionnaires déjà étudiées vont nous permettre
d'interroger une \textbf{table de données} à l'aide d'opération appelées
\textbf{requêtes} :

\begin{itemize}
\tightlist
\item
  pour tester l'appartenance à la table d'un enregistrement vérifiant
  une certaine condition sur ses attributs ;
\item
  pour calculer une valeur en combinant les valeurs d'attributs de
  plusieurs enregistrements (\textbf{opération d'agrégation}) ;
\item
  pour sélectionner tous les enregistrements vérifiant une certaine
  condition sur leurs attributs (\textbf{opération de sélection par
  lignes}) ;
\item
  pour projeter les enregistrements sur certains de leurs attributs
  (\textbf{opération de projection sur des colonnes}) ;
\item
  pour trier les enregistrements en fontion d'un ordre prédéfini sur
  leurs attributs (\textbf{opération de tri}).
\end{itemize}

\end{definition}

\hypertarget{requuxeates-de-recherche}{%
\section{Requêtes de recherche}\label{requuxeates-de-recherche}}

\begin{exercice}{}

\begin{enumerate}
\def\labelenumi{\arabic{enumi}.}
\tightlist
\item
  Ouvrir le script
  \passthrough{\lstinline!'TP\_Recherche\_Tris\_Eleve.py'!} dans un IDE
  Python
\item
  Charger dans une variable \passthrough{\lstinline!table\_clients!} la
  table contenue dans le fichier \passthrough{\lstinline!'clients.csv'!}
  à l'aide de la fonction \passthrough{\lstinline!lecture\_csv!} qui est
  fournie. Il s'agit de la table des clients d'un site marchand que nous
  avons manipulée dans le cours. Exécuter le test unitaire
  \passthrough{\lstinline!test\_import\_table\_clients(table\_clients)!}.
\end{enumerate}

\begin{lstlisting}[language=Python]
table_clients = lecture_csv('clients.csv', ',')
\end{lstlisting}

\begin{enumerate}
\def\labelenumi{\arabic{enumi}.}
\setcounter{enumi}{2}
\tightlist
\item
  Compléter la fonction \passthrough{\lstinline!recherche\_attribut!}
  ci-dessous en respectant sa spécification. Vérifier en exécutant le
  test unitaire
  \passthrough{\lstinline!test\_recherche\_attribut\_table\_clients(table\_clients)!}.
\end{enumerate}

\begin{lstlisting}[language=Python]
def recherche_attribut(table, attribut, valeur):
    """Paramètres : 
        table un tableau de dictionnaires, table  de clients.csv
        attribut de type str  valeur du type d'attribut dans table
    Valeur renvoyée:
        Un booléen indiquant si table contient un enregistrement e
        tel que e[attribut] == valeur"""
    for enregistrement in table:
        if .........................:
            return ............
    return .........
\end{lstlisting}

\begin{enumerate}
\def\labelenumi{\arabic{enumi}.}
\setcounter{enumi}{3}
\tightlist
\item
  Compléter la fonction
  \passthrough{\lstinline!recherche\_attribut\_et!} ci-dessous en
  respectant sa spécification. Exécuter le test unitaire
  \passthrough{\lstinline!test\_recherche\_attributs\_et(table\_clients)!}.
\end{enumerate}

\begin{lstlisting}[language=Python]
def recherche_attributs_et(table, attribut1, valeur1, attribut2, valeur2):
    """Paramètres : 
        table un tableau de dictionnaires, table  de clients.csv
        attribut1 de type str, valeur1 du type d'attribut1
        attribut2 de type str, valeur2 du type d'attribut2
    Valeur renvoyée:
        Un booléen indiquant si table contient un enregistrement e
        tel que e[attribut1] == valeur1 et e[attribut2] == valeur2 """
    for enregistrement in table:
        if .........................:
            return ............
    return .........
\end{lstlisting}

Proposer une modification de
\passthrough{\lstinline!recherche\_attributs\_et!} qui prend en
paramètres deux tableaux \passthrough{\lstinline!attribut!} et
\passthrough{\lstinline!valeur!} et recherche si la table contient un
enregistrement pour lequel tous les attributs listés correspondent aux
valeurs listées puis un enregistrement pour lequel au moins un des
attributs listés correspond à son homologue dans la liste des valeurs.

\end{exercice}

\hypertarget{requuxeates-avec-opuxe9ration-dagruxe9gation}{%
\section{Requêtes avec opération
d'agrégation}\label{requuxeates-avec-opuxe9ration-dagruxe9gation}}

\begin{exercice}{}

On travaille toujours avec la table contenue dans le fichier
\passthrough{\lstinline!'clients.csv'!} avec le script
\passthrough{\lstinline!'TP\_Recherche\_Tris\_Eleve.py'!} dans un IDE
Python.

\begin{enumerate}
\def\labelenumi{\arabic{enumi}.}
\tightlist
\item
  Compléter la fonction \passthrough{\lstinline!nombre\_departement!}
  ci-dessous en respectant sa spécification. On donne une postcondition
  qui doit être vérifiée.
\end{enumerate}

\begin{lstlisting}[language=Python]
def nombre_departement(table, departement):
    """Paramètres : 
        table un tableau de dictionnaires, table  de clients.csv
        departement de type str, un numéro de département
    Valeur renvoyée :
        Nombre d'occurences de departement dans table"""
    compteur = 0
    for enregistrement in table:
        if ....................................:
            compteur = .........................
    return compteur

# postcondition
assert nombre_departement(table_clients, "département", "69") == 481
\end{lstlisting}

\begin{enumerate}
\def\labelenumi{\arabic{enumi}.}
\setcounter{enumi}{1}
\tightlist
\item
  Compléter la fonction \passthrough{\lstinline!nombre\_departement!}
  ci-dessous en respectant sa spécification. Exécuter le test unitaire
  \passthrough{\lstinline!test\_nombre\_occurences(table\_clients)!}.
\end{enumerate}

\begin{lstlisting}[language=Python]
def nombre_occurences(table, attribut, valeur):
    """Paramètres : 
        table un tableau de dictionnaires, table  de clients.csv
        attribut de type str, valeur du type d'attribut dans table
    Valeur renvoyée :
        Nombre d'occurences d'attribut avec valeur dans table"""
    ...........................
    ...........................
    ...........................
    ...........................
    ...........................
    ...........................
\end{lstlisting}

\begin{enumerate}
\def\labelenumi{\arabic{enumi}.}
\setcounter{enumi}{2}
\tightlist
\item
  Compléter la fonction \passthrough{\lstinline!moyenne\_visites!}
  ci-dessous en respectant sa spécification. On donne une postcondition
  qui doit être vérifiée.
\end{enumerate}

\begin{lstlisting}[language=Python]
def moyenne_visites(table):
    """Paramètres : 
        table un tableau de dictionnaires, table  de clients.csv
    Valeur renvoyée :
        Moyenne des visites par enregistrement de type float"""
    somme = 0
    taille = 0
    for enregistrement in table:
        .......................
        .......................
        .......................
    return somme / taille

#postcondition
assert moyenne_visites(table_clients) == 76.2807 
\end{lstlisting}

\begin{enumerate}
\def\labelenumi{\arabic{enumi}.}
\setcounter{enumi}{3}
\tightlist
\item
  Compléter la fonction \passthrough{\lstinline!minimum\_visites!}
  ci-dessous en respectant sa spécification. On donne une postcondition
  qui doit être vérifiée.
\end{enumerate}

\begin{lstlisting}[language=Python]
def minimum_visites(table):
    """Paramètres : 
        table un tableau de dictionnaires, table  de clients.csv
    Valeur renvoyée :
        nombre minimum de visites de type int"""    
    min_visites = int(table[0]['visites'])
    for enregistrement in table[1:]:
        ............................
        ............................
        ............................
    return min_visites

#postcondition
assert minimum_visites(table_clients)  == 2
\end{lstlisting}

\begin{enumerate}
\def\labelenumi{\arabic{enumi}.}
\setcounter{enumi}{4}
\tightlist
\item
  Compléter la fonction
  \passthrough{\lstinline!departement\_max\_occurence!} ci-dessous en
  respectant sa spécification. On donne une postcondition qui doit être
  vérifiée.
\end{enumerate}

\begin{lstlisting}[language=Python]
def departement_max_occurence(table):
    """Paramètre : table sous forme de tableau de dictionnaires
    Valeur renvoyée : tuple formé du nombre d'occurences maximal parmi les départements
    et du tableau des  départements réalisant ce maximum
    """
    ............
    ............ 













assert departement_max_occurence(table_clients) == (547, ['59'])
\end{lstlisting}

\end{exercice}

\hypertarget{requuxeates-avec-suxe9lection-de-lignes}{%
\section{Requêtes avec sélection de
lignes}\label{requuxeates-avec-suxe9lection-de-lignes}}

\begin{exercice}{}

On travaille toujours avec la table contenue dans le fichier
\passthrough{\lstinline!'clients.csv'!} avec le script
\passthrough{\lstinline!'TP\_Recherche\_Tris\_Eleve.py'!} dans un IDE
Python.

\emph{Un opération de \textbf{sélection} consiste à construire une
nouvelle table avec les mêmes attributs mais en filtrant les
enregistrements (ou lignes) selon une condition logique.}

\begin{enumerate}
\def\labelenumi{\arabic{enumi}.}
\tightlist
\item
  Compléter la fonction \passthrough{\lstinline!selection\_departement!}
  ci-dessous en respectant sa spécification. On donne une postcondition
  qui doit être vérifiée.
\end{enumerate}

\begin{lstlisting}[language=Python]
def selection_departement(table, departement):
    """
    Paramètres : 
        table une table sous forme de tableau de dictionnaires
        department une chaine de caractères représentant un département
    Valeur renvoyée :
        tableau de dictionnaires contenant les enregistrement de table dont l'attribut "département"  a la valeur passée en paramètre
    """
    return [enregistrement
             for enregistrement in table if ......................................    ]

# postcondition
assert selection_departement(table_clients, "69")[0]['email'] == 'nnguyen@noos.fr' \
        and len(selection_departement(table_clients, "69")) == 481
\end{lstlisting}

\begin{enumerate}
\def\labelenumi{\arabic{enumi}.}
\setcounter{enumi}{1}
\tightlist
\item
  Compléter la fonction
  \passthrough{\lstinline!selection\_depart\_visites\_min!} ci-dessous
  en respectant sa spécification. On donne une postcondition qui doit
  être vérifiée.
\end{enumerate}

\begin{lstlisting}[language=Python]
def selection_depart_visites_min(table, departement, visites_min):
    """Paramètres : 
        table une table sous forme de tableau de dictionnaires
        department une chaine de caractères représentant un département
        visites_min un entier naturel 
    Valeur renvoyée :
        tableau de dictionnaires contenant les enregistrement de table dont l'attribut "département"  a la valeur passée en paramètre et l'attribut visites est >= visites_min"""
    return [enregistrement for enregistrement in table 
            if ........................................................]         


assert len(selection_depart_visites_min(table_clients,"69",100))==171
\end{lstlisting}

\end{exercice}

\hypertarget{requuxeates-avec-projection-sur-des-colonnes}{%
\section{Requêtes avec projection sur des
colonnes}\label{requuxeates-avec-projection-sur-des-colonnes}}

\begin{exercice}{}

On travaille toujours avec la table contenue dans le fichier
\passthrough{\lstinline!'clients.csv'!} avec le script
\passthrough{\lstinline!'TP\_Recherche\_Tris\_Eleve.py'!} dans un IDE
Python.

\emph{Une opération de \textbf{projection} consiste à construire une
nouvelle table en filtrant les attributs.}

\begin{enumerate}
\def\labelenumi{\arabic{enumi}.}
\tightlist
\item
  Compléter la fonction \passthrough{\lstinline!projection\_visites!}
  ci-dessous en respectant sa spécification. On donne une postcondition
  qui doit être vérifiée.
\end{enumerate}

\begin{lstlisting}[language=Python]
def projection_visites(table):
    """Paramètres : 
        table une table sous forme de tableau de dictionnaires
    Valeur renvoyée :
        tableau des valeurs des attributs "visites" pour les enregistrements de table
        avec conversion des nombres de visites en entiers
    """
    return [ ......................... for enregistrement in table ]

assert projection_visites(table_clients)[:10] == [57, 145, 67, 131, 76, 52, 65, 3, 101, 18]
\end{lstlisting}

\begin{enumerate}
\def\labelenumi{\arabic{enumi}.}
\setcounter{enumi}{1}
\tightlist
\item
  Compléter la fonction \passthrough{\lstinline!projection\_visites!}
  ci-dessous en respectant sa spécification. On donne une postcondition
  qui doit être vérifiée.
\end{enumerate}

\begin{lstlisting}[language=Python]
def selection_departement_projection_visites(table, departement):
    """Paramètres : 
        table une table sous forme de tableau de dictionnaires
    Valeur renvoyée :
        tableau des valeurs des attributs "visites" pour les enregistrements de table du département
        passé en paramètre avec conversion des nombres de visites en entiers"""
    ...................................................................
    ...................................................................

assert selection_departement_projection_visites(table_clients, "69")[:10] == [43, 52, 127, 53, 41, 117, 31, 86, 107, 145]
\end{lstlisting}

\begin{enumerate}
\def\labelenumi{\arabic{enumi}.}
\setcounter{enumi}{2}
\tightlist
\item
  Compléter la fonction
  \passthrough{\lstinline!projection\_departement\_age!} ci-dessous en
  respectant sa spécification. On donne une postcondition qui doit être
  vérifiée.
\end{enumerate}

\begin{lstlisting}[language=Python]
def projection_departement_age(table, annee):
    """Paramètres : 
        table une table sous forme de tableau de dictionnaires
    Valeur renvoyée :
        tableau de dictionnaires avec deux attributs le département et l'âge du client
        calculé à partir de sa date de naissance."""
    return [{'département' : enreg['département'], 
            'âge' : ................................................} 
           for enreg in table]

assert projection_departement_age(table_clients, 2021)[:3] == [{'département': '79', 'âge': 61}, {'département': '10', 'âge': 27},
 {'département': '73', 'âge': 20}]
\end{lstlisting}

\end{exercice}

\hypertarget{tri-dune-table-de-donnuxe9es}{%
\section{Tri d'une table de
données}\label{tri-dune-table-de-donnuxe9es}}

\begin{methode}{}

Pour trier une \textbf{table de données} implémentée sous forme de
\textbf{tableau de dictionnaires} en Python, il est préférable
d'utiliser les fonctions de tri de la bibliothèque standard plutôt que
des tris que nous avons programmés.

D'une part nous avons ainsi la garantie d'une complexité optimale en
\(O(nlog(n))\), d'autre part les fonctions de tri de bibliothèque
permettent de paramétrer le critère de tri appelé aussi \textbf{clef de
tri}.

Il existe deux fonctions de tri dans la bibliothèque standard :

\begin{itemize}
\item
  \passthrough{\lstinline!sorted!} s'applique à un tableau
  \passthrough{\lstinline!tab!} avec la syntaxe
  \passthrough{\lstinline!sorted(tab, key = clef\_tri)!} et renvoie un
  nouveau tableau avec les éléments du tableau
  \passthrough{\lstinline!tab!} triés selon la clef de tri passée en
  paramètre. Celle-ci est une fonction qui s'applique à un élément du
  tableau et renvoie une valeur calculée à partir de celle de l'élément.
  Le tri s'effectue en comparant les valeurs renvoyées par la clef de
  tri pour chaque élément. Cette valeur peut être un
  \passthrough{\lstinline!tuple!} puisque dans une table, un élément est
  un dictionnaire avec éventuellement plusieurs attributs. Dans ce cas,
  pour comparer deux valeurs, on applique \textbf{l'ordre
  lexicographique} : on compare les premières composantes, puis les
  secondes etc \ldots. Par défaut la comparaison s'effectue dans l'ordre
  croissant, mais avec le paramètre optionnel
  \passthrough{\lstinline!reverse!} on peut demander un ordre
  décroissant :
  \passthrough{\lstinline!sorted(tab, key = clef\_tri, reverse = True)!}.
\item
  \passthrough{\lstinline!sort!} est une fonction de tri en place, elle
  ne renvoie pas un nouveau tableau. Elle s'applique à un tableau
  \passthrough{\lstinline!tab!} avec la syntaxe
  \passthrough{\lstinline!tab.sort(key = clef\_tri)!}, les paramètres
  sont les mêmes que pour \passthrough{\lstinline!sorted!}.
\end{itemize}

Attention, \passthrough{\lstinline!sorted!} ne renvoie qu'une copie
superficielle du tableau ! Si ses éléments sont des références et c'est
le cas de l'implémentation des \textbf{tables de données} sous forme de
\textbf{tableaux de dictionnaires}, il faut effectuer une copie profonde
du tableau avec la fonction \passthrough{\lstinline!deepcopy!} du module
\passthrough{\lstinline!copy!} pour obtenir une vraie copie triée du
tableau initiale :
\passthrough{\lstinline!sorted(deepcopy(tab), key = clef\_tri)!}.

On donne ci-dessous quelques exemples. Notons que pour obtenir un tri
décroissant selon l'attribut `\passthrough{\lstinline!note!}' puis
croissant selon l'attribut \passthrough{\lstinline!'langage'!}, on ne
peut pas utiliser une clef de tri qui renvoie le couple d'attributs
(\passthrough{\lstinline!'note'!}, \passthrough{\lstinline!'langage'!})
car les ordres ne sont pas les mêmes selon les composantes. On procède
par composition des tris avec \passthrough{\lstinline!sorted!} dans
l'ordre inverse des priorités de tri : d'abord selon l'attribut
\passthrough{\lstinline!note!} puis selon l'attribut
\passthrough{\lstinline!langage!}.

Une propriété importante des fonctions de tri de bibliothèque
\passthrough{\lstinline!sorted!} et \passthrough{\lstinline!sort!} est
la \textbf{stabilité du tri} : si on enchaîne deux tris successifs (par
ordre lexicographique ou composition de
\passthrough{\lstinline!sorted!}), deux éléments égaux pour le second
tri conservent l'ordre du premier tri.

\begin{lstlisting}[language=Python]
>>> table = [{'élève' : 'guido', 'langage' : 'python', 'note' : 19}, {'élève' : 'monty', 'langage' : 'python', 'note' : 20}, {'élève' : 'brian', 'langage' : 'c', 'note' : 20}]
>>> def clef_note(enreg):
...     return enreg['note']
>>> def clef_langage(enreg):
...     return enreg['langage']
>>> def clef_langage_note(enreg):
...     return (enreg['langage'], enreg['note'])
>>> def clef_note_langage(enreg):
...     return (enreg['note'], enreg['langage'])
>>> sorted(table, key = clef_note)
[{'élève': 'guido', 'langage': 'python', 'note': 19}, {'élève': 'monty', 'langage': 'python', 'note': 20}, {'élève': 'brian', 'langage': 'c', 'note': 20}]
>>> sorted(table, key = clef_note, reverse = True)
[{'élève': 'monty', 'langage': 'python', 'note': 20}, {'élève': 'brian', 'langage': 'c', 'note': 20}, {'élève': 'guido', 'langage': 'python', 'note': 19}]
>>> sorted(table, key = clef_langage_note)  #ordre lexicographique (langage croissant, note croissant)
[{'élève': 'brian', 'langage': 'c', 'note': 20}, {'élève': 'guido', 'langage': 'python', 'note': 19}, {'élève': 'monty', 'langage': 'python', 'note': 20}]
>>> sorted(table, key = clef_note_langage)  #ordre lexicographique (note croissant, langage croissant)
[{'élève': 'guido', 'langage': 'python', 'note': 19}, {'élève': 'brian', 'langage': 'c', 'note': 20}, {'élève': 'monty', 'langage': 'python', 'note': 20}]
>>> table_tri_lang_cr = sorted(table, key = clef_langage)   #on va composer les tris avec sorted
>>> sorted(table_tri_lang_cr, key = clef_note, reverse = True) # équivalent à (note décroissant, langage croissant)
[{'élève': 'brian', 'langage': 'c', 'note': 20}, {'élève': 'monty', 'langage': 'python', 'note': 20}, {'élève': 'guido', 'langage': 'python', 'note': 19}]
\end{lstlisting}

\end{methode}

\begin{exercice}{}

On travaille toujours avec la table contenue dans le fichier
\passthrough{\lstinline!'clients.csv'!} avec le script
\passthrough{\lstinline!'TP\_Recherche\_Tris\_Eleve.py'!} dans un IDE
Python.

\begin{enumerate}
\def\labelenumi{\arabic{enumi}.}
\tightlist
\item
  Compléter la clef de tri \passthrough{\lstinline!clef\_departement!}
  pour que \passthrough{\lstinline!table\_tri\_departement!} soit trié
  dans l'ordre croissant des numéros de département dans le code
  ci-dessous. Vérifier la postcondition donnée dans le fichier.
\end{enumerate}

\begin{lstlisting}[language=Python]
def clef_departement(enreg):
    return .........................

table_tri_departement = sorted(table_clients, key = clef_departement)
\end{lstlisting}

\begin{enumerate}
\def\labelenumi{\arabic{enumi}.}
\setcounter{enumi}{1}
\tightlist
\item
  Compléter la clef de tri \passthrough{\lstinline!clef\_visites!} pour
  que \passthrough{\lstinline!table\_tri\_visites\_decroissant!} soit
  trié dans l'ordre décroissant des nombres de visites dans le code
  ci-dessous. Vérifier la postcondition donnée dans le fichier.
\end{enumerate}

\begin{lstlisting}[language=Python]
def clef_visites(enreg):
    return .........................

table_tri_visites_decroissant = sorted(table_clients, key = clef_visites, reverse = True)
\end{lstlisting}

\begin{enumerate}
\def\labelenumi{\arabic{enumi}.}
\setcounter{enumi}{2}
\tightlist
\item
  Compléter la clef de tri
  \passthrough{\lstinline!clef\_departement\_visites!} pour que
  \passthrough{\lstinline!table\_tri\_dep\_vis\_croissant!} soit trié
  dans l'ordre croissant des départements puis des nombres de visites
  dans le code ci-dessous. Vérifier la postcondition donnée dans le
  fichier.
\end{enumerate}

\begin{lstlisting}[language=Python]
def clef_departement_visites(enreg):
    return .........................

table_tri_dep_vis_croissant = sorted(table_clients, key = clef_departement_visites)
\end{lstlisting}

\begin{enumerate}
\def\labelenumi{\arabic{enumi}.}
\setcounter{enumi}{3}
\tightlist
\item
  Proposer une instruction qui permette de trier
  \passthrough{\lstinline!table\_clients!} d'abord par département
  croissant puis par nombre de visites décroissant. Vérifier la
  postcondition donnée dans le fichier.
\end{enumerate}

\begin{lstlisting}[language=Python]
table_tri_dep_crois_vis_decrois =  ...............
                                   ...............
\end{lstlisting}

\begin{enumerate}
\def\labelenumi{\arabic{enumi}.}
\setcounter{enumi}{4}
\tightlist
\item
  Le fichier \passthrough{\lstinline!transactions.csv!} contient une
  table de données de nouvelles transactions effectuées sur le site
  marchand. Chaque transaction est identifiée par deux attributs :
  \passthrough{\lstinline!'email'!} pour l'email du client et
  '\passthrough{\lstinline!dépenses'!} pour le montant de la dépense.
\end{enumerate}

\begin{lstlisting}
email,dépenses
wpereira@orange.fr,104.91
ariviere@tiscali.fr,18.37
\end{lstlisting}

On souhaite mettre à jour les attributs
\passthrough{\lstinline!'visites'!} et
\passthrough{\lstinline!'dépenses'!} de
\passthrough{\lstinline!table\_clients!} avec les nouvelles
transactions. On considère que chaque client est identifié de façon
unique par son email et on peut donc insérer efficacement les valeurs de
chaque transaction dans \passthrough{\lstinline!table\_clients!} avec
une recherche dichotomique sur l'attribut
\passthrough{\lstinline!'email'!}. Bien sûr, il faut d'abord trier
\passthrough{\lstinline!table\_clients!} selon l'attribut
\passthrough{\lstinline!'email'!}. Compléter le code ci-dessous.
Décommenter les postconditions dans le fichier pour vérifier le code et
contrôler le contenu du fichier de sortie
\passthrough{\lstinline!'clients\_maj.csv'!}.

\begin{lstlisting}[language=Python]
from copy import deepcopy #pour réaliser une copie de table

def recherche_dicho_croissant(element, table, attribut):
    """Paramètres :
         table un tableau de dictionnaires
         attribut de type str
         element une valeur possible pour l'attribut
    Valeur renovoyée : index de la valeur element de attribut dans table"""
    debut = 0
    fin = len(table) - 1
    while fin - debut >= 0:
        milieu = (debut + fin) // 2
        if table[milieu][attribut] < element:
            debut = .....................
        elif table[milieu][attribut] > element:
            fin = .......................
        else:
            return ......................
    return None

def clef_email(enreg):
    return enreg['email']

def maj_depenses_table(table, transactions): 
    """Paramètres : table et transactions deux  tables sous forme de tableaux de dictionnaires
    Valeur renvoyée :
        table_tri un tableau de dictionnaires mise à jour des attributs 'visites' et 'dépenses'
        de table par les valeurs de transactions""""       
    table_tri = sorted(deepcopy(table), key = clef_email)
    table_cible = []
    for enreg in transactions:
        index_email = recherche_dicho_croissant(enreg['email'], table_tri, 'email')
        if index_email is not None:
            ..........................................
            ..........................................
            ..........................................
    return table_tri
                     
table_clients = lecture_csv('clients.csv', ',')
transactions = lecture_csv('transactions.csv', ',') 
table_tri = maj_depenses_table(table_clients, transactions)
ecriture_csv(table_tri, 'clients_maj.csv', ',')
\end{lstlisting}

\end{exercice}

\hypertarget{applications-uxe0-un-autre-exemple}{%
\section{Applications à un autre
exemple}\label{applications-uxe0-un-autre-exemple}}

\begin{exercice}{}

Nous allons utiliser un fichier nommé
\passthrough{\lstinline!countries.csv!} qui contient quelques données
sur les différents pays du monde. En voici les premières lignes : les
champs sont clairement séparés par des points-virgules.

\begin{lstlisting}
iso;name;area;population;continent;currency_code;currency_name;capital
AD;Andorra;468.0;84000;EU;EUR;Euro;6
AE;United Arab Emirates;82880.0;4975593;AS;AED;Dirham;21
AF;Afghanistan;647500.0;29121286;AS;AFN;Afghani;81
\end{lstlisting}

Les données sont issues du site \url{http://www.geonames.org} et ont été
légèrement simplifiées. La signification des différents champs est
transparente (\passthrough{\lstinline!currency!} signifie devise), à
part le dernier champ, nommé capital et dont les valeurs sont des
numéros d'identifiants de villes que l'on trouvera dans un autre fichier
nommé \passthrough{\lstinline!cities.csv!} que nous utiliserons dans le
chapitre sur les fusions de tables.

\begin{enumerate}
\def\labelenumi{\arabic{enumi}.}
\tightlist
\item
  On travaille toujours avec le script
  \passthrough{\lstinline!'TP\_Recherche\_Tris\_Eleve.py'!} dans un IDE
  Python à la suite des exercices précédents. On commence par charger la
  table avec la fonction \passthrough{\lstinline!lecture\_csv!},
  attention le délimiteur de champ n'est pas le symbole
  \passthrough{\lstinline!,!} mais \passthrough{\lstinline!;!}.
\end{enumerate}

\begin{lstlisting}[language=Python]
table_pays = lecture_csv('countries.csv', ';')
\end{lstlisting}

\begin{enumerate}
\def\labelenumi{\arabic{enumi}.}
\setcounter{enumi}{1}
\tightlist
\item
  Compléter la fonction \passthrough{\lstinline!nombre\_europe!}
  ci-dessous en respectant sa spécification. On donne une postcondition
  qui doit être vérifiée.
\end{enumerate}

\begin{lstlisting}[language=Python]
def nombre_europe(table):
    """Paramètre : table de countries.csv
    Valeur renvoyée : compteur de type int représentant le nombre de pays du continent européen
    """
    compteur = 0
    for enregistrement in table:
        .......................
        ......................
    return compteur

assert nombre_europe(table_pays) == 52
\end{lstlisting}

\begin{enumerate}
\def\labelenumi{\arabic{enumi}.}
\setcounter{enumi}{2}
\tightlist
\item
  Compléter la fonction \passthrough{\lstinline!selection\_europe!}
  ci-dessous en respectant sa spécification. On donne une postcondition
  qui doit être vérifiée.
\end{enumerate}

\begin{lstlisting}[language=Python]
def selection_europe(table):
    """Paramètre : table de countries.csv
    Valeur renvoyée : table des enregistrements des pays du continent européen"""
    return ................................


europe = selection_europe(table_pays)
assert len(europe) == 52 and europe[0]['name'] == 'Andorra'
\end{lstlisting}

\begin{enumerate}
\def\labelenumi{\arabic{enumi}.}
\setcounter{enumi}{3}
\tightlist
\item
  Compléter la fonction
  \passthrough{\lstinline!selection\_europe\_non\_euro!} ci-dessous en
  respectant sa spécification. On donne une postcondition qui doit être
  vérifiée.
\end{enumerate}

\begin{lstlisting}[language=Python]
def selection_europe_non_euro(table):
     """Paramètre : table de countries.csv
    Valeur renvoyée : table des enregistrements des pays du continent européen qui n'ont pas pour monnaie l'euro"""
    return ................................
           ................................

europe_non_euro = selection_europe_non_euro(table_pays)
assert len(europe_non_euro) == 27 and europe_non_euro[0]['name'] == 'Albania'
\end{lstlisting}

\begin{enumerate}
\def\labelenumi{\arabic{enumi}.}
\setcounter{enumi}{4}
\tightlist
\item
  Compléter la fonction \passthrough{\lstinline!projection\_aire!}
  ci-dessous en respectant sa spécification. On donne une postcondition
  qui doit être vérifiée.
\end{enumerate}

\begin{lstlisting}[language=Python]
def projection_aire(table):
    """Paramètre : table de countries.csv
    Valeur renvoyée : tableau des aires (type float) de tous les enregistrements"""
    return ................................

assert  projection_aire(table_pays)[:5]  == [468.0, 82880.0, 647500.0, 443.0, 102.0]
\end{lstlisting}

\begin{enumerate}
\def\labelenumi{\arabic{enumi}.}
\setcounter{enumi}{5}
\tightlist
\item
  Compléter la fonction
  \passthrough{\lstinline!projection\_pays\_densite!} ci-dessous en
  respectant sa spécification. On donne une postcondition qui doit être
  vérifiée.
\end{enumerate}

\begin{lstlisting}[language=Python]
def projection_pays_densite(table):
    """Paramètre : table de countries.csv
    Valeur renvoyée : nouvelle table avec deux attributs 'pays' et 'densité' de population"""
    return [ {'pays' : enreg['name'], 
             'densité' : ...............................)}
            for enreg in table]

assert projection_pays_densite(table_pays)[:3] == [{'pays': 'Andorra', 'densité': 179.48717948717947},
 {'pays': 'United Arab Emirates', 'densité': 60.033699324324324},
 {'pays': 'Afghanistan', 'densité': 44.974959073359074}]
\end{lstlisting}

\begin{enumerate}
\def\labelenumi{\arabic{enumi}.}
\setcounter{enumi}{6}
\tightlist
\item
  Écrire une fonction \passthrough{\lstinline!maximum\_densite!}
  respectant la spécification ci-dessous. On donne une postcondition qui
  doit être vérifiée.
\end{enumerate}

\begin{lstlisting}[language=Python]
def maximum_densite(table):
    """Paramètre : table de countries.csv
    Valeur renvoyée : tuple avec le nom du pays de densité maximale de population et cette densité maximale"""











assert maximum_densite(table_pays) == ('Monaco', 16905.128205128207)
\end{lstlisting}

\begin{enumerate}
\def\labelenumi{\arabic{enumi}.}
\setcounter{enumi}{7}
\tightlist
\item
  Écrire une fonction
  \passthrough{\lstinline!population\_par\_continent!} respectant la
  spécification ci-dessous. On donne une postcondition qui doit être
  vérifiée.
\end{enumerate}

\begin{lstlisting}[language=Python]
def population_par_continent(table):
    """Paramètre : table de countries.csv
    Valeur renvoyée : dictionnaire de clefs les identifiants des continents
    et de valeurs les populations cumulées des pays leur appartenant"""
    








        

assert population_par_continent(table_pays) == {'EU': 740017414,
 'AS': 4119426856,  'NA': 539886359, 'AF': 1018849428, 'SA': 400143568, 
 'OC': 36066083}
\end{lstlisting}

\begin{enumerate}
\def\labelenumi{\arabic{enumi}.}
\setcounter{enumi}{8}
\tightlist
\item
  Écrire une fonction \passthrough{\lstinline!densite\_max\_top5!}
  respectant la spécification ci-dessous. On donne une postcondition qui
  doit être vérifiée.
\end{enumerate}

\begin{lstlisting}[language=Python]
def densite_max_top5(table):
     """Paramètre : table de countries.csv
    Valeur renvoyée : table avec les noms et les densités des 5 pays les plus  densément peuplées dans l'ordre décroissant des densités de population"""
    












assert densite_max_top5(table_pays) == [{'pays': 'Monaco', 'densité': 16905.128205128207},
 {'pays': 'Singapore', 'densité': 6786.5872672152445},
 {'pays': 'Hong Kong', 'densité': 6317.478021978022},
 {'pays': 'Gibraltar', 'densité': 4289.846153846154},
 {'pays': 'Vatican', 'densité': 2093.181818181818}]
\end{lstlisting}

\end{exercice}

\hypertarget{synthuxe8se}{%
\section{Synthèse}\label{synthuxe8se}}

\begin{synthese}{}

Lorsqu'une \textbf{table de données} contenue dans un \textbf{fichier
CSV} est chargée dans une structure de données Python comme un
\textbf{tableau de dictionnaires}, on peut la manipuler avec des
requêtes de \textbf{recherche}, \textbf{d'agrégation}, de
\textbf{sélection sur les lignes} ou de \textbf{projection sur les
colonnes}. Il est possible également de trier les enregistrements d'une
table avec la fonction de bibliothèque \passthrough{\lstinline!sorted!}
en lui passant une fonction clef de tri. Elle garantit la
\textbf{stabilité du tri} : les éléments égaux conservent leur ordre
intial. On peut ainsi extraire des informations d'une table ou
construire de nouvelles tables. En classe de terminale, nous étudierons
les \textbf{bases de données} dans le modèle relationnel, qui peuvent
être modélisées par des tables. Les requêtes pour les interroger seront
similaires mais exprimées dans un langage spécifique, le
\href{https://pixees.fr/informatiquelycee/n_site/nsi_term_bd_sql.html}{SQL}.

\end{synthese}

\end{document}
